\section{Empirical Temporal Functionals}
\label{app:functionals}

The controlled experiments separate the dynamics that generate the inputs from
the temporal map that generates the response.  This appendix defines the task
family used throughout the synthetic study and fixes its normalization,
delays, and allocation across experiments.  The hierarchy follows the usual
reservoir-computing distinction between linear memory and nonlinear
information processing \citep{Jaeger2001,DambreEtAl2012}, while the exponentially
weighted and Volterra-type tasks are finite-window versions of classical
fading-memory operators \citep{BoydChua1985}.

\subsection{Notation and common normalization}
\label{app:functional-notation}

Let
\[
    \bm X_{t,w}:=(X_{t-w+1},\ldots,X_t)\in([-1,1]^d)^w
\]
be the input window presented to the learner.  Choose directions
\(u_d,v_d\in\R^d\) satisfying
\(\norm{u_d}_1,\norm{v_d}_1\leq1\), and fix the memory-decay parameter
\(\gamma=0.8\).  For the base dimension \(d=3\), the common directions are
\[
    u_3=\frac13(1,-1,1)^\top,
    \qquad
    v_3=\frac13(1,1,-1)^\top.
\]
All task directions, delays, decay parameters, and any task-generation seeds
are fixed before the data are simulated; all methods evaluated on a given
dataset receive the same realized task. Define
\begin{equation}
    A_w:=\sum_{j=0}^{w-1}\gamma^j
        =\frac{1-\gamma^w}{1-\gamma}.
    \label{eq:functional-Aw}
\end{equation}
The normalizations below place every controlled target in \([-1,1]\), so the
bounded-label constant can be taken as \(\Upsilon=1\) for this suite.

\subsection{Shared temporal tasks}
\label{app:shared-functionals}

\paragraph{One-step future projection.}
The genuine forecasting task uses one observation beyond the input window,
\begin{equation}
    F_{\mathrm{fore}}(\bm X_{t,w},X_{t+1};u_d)
    :=u_d^\top X_{t+1}.
    \label{eq:functional-forecast}
\end{equation}
The learner receives observations only through time \(t\); \(X_{t+1}\) is used
solely to construct the label.  This distinguishes the task from an
instantaneous decoding check based on \(X_t\).

\paragraph{Exponentially weighted linear memory.}
Let
\begin{equation}
    L_u(t)
    :=\sum_{j=0}^{w-1}\gamma^j u_d^\top X_{t-j}.
    \label{eq:functional-linear-memory}
\end{equation}
The normalized fading-memory target is
\begin{equation}
    F_{\mathrm{exp}}(\bm X_{t,w};u_d)
    :=\frac{L_u(t)}{A_w}.
    \label{eq:functional-exp}
\end{equation}
It requires information from the full window while assigning greater weight to
recent observations.

\paragraph{Second-order Volterra-type memory.}
With \(L_v(t)\) defined as in \cref{eq:functional-linear-memory} using \(v_d\),
set
\begin{equation}
    F_{\mathrm{vol}}(\bm X_{t,w};u_d,v_d)
    :=
    \frac{L_u(t)+\tfrac12 L_v(t)^2}
         {A_w+\tfrac12 A_w^2}.
    \label{eq:functional-volterra}
\end{equation}
This task combines distributed linear memory with a smooth nonlinear
interaction across lags.

\paragraph{Delayed recall.}
For a delay \(\tau\in\{1,5,10,20\}\subset\{0,\ldots,w-1\}\), define
\begin{equation}
    F_{\mathrm{delay},\tau}(\bm X_{t,w};u_d)
    :=u_d^\top X_{t-\tau}.
    \label{eq:functional-delay}
\end{equation}
The resulting delay curve is a direct diagnostic of how much information about
a particular past observation remains in the final reservoir feature vector.
It is interpreted as an empirical memory profile, not as an estimator of the
worst-case contraction bound.

\paragraph{Sparse cross-lag interaction.}
For the predeclared delays \((\tau_1,\tau_2)=(0,15)\), define
\begin{equation}
    F_{\mathrm{cross}}(\bm X_{t,w};u_d,v_d)
    :=
    \bigl(u_d^\top X_{t-\tau_1}\bigr)
    \bigl(v_d^\top X_{t-\tau_2}\bigr).
    \label{eq:functional-cross}
\end{equation}
Unlike the dense Volterra target, this functional isolates one nonlinear
interaction between a recent and a distant observation.

\subsection{Family-specific tasks}
\label{app:family-specific-functionals}

The common tasks make performance comparable across generators.  Two
additional targets exploit the defining latent structure of the HMM and GARCH
families.

\paragraph{HMM filtered regime score.}
For the three-state HMM in \cref{app:hmm-family}, assign state scores
\(a=(-1,0,1)^\top\).  Using the known generator parameters, compute the
finite-window filtering probabilities
\begin{equation}
    \pi_{t,w}(j)
    :=\Prob\!\left(
      S_t=j\mid X_{t-w+1:t}
    \right),
    \qquad j\in\{1,2,3\},
    \label{eq:functional-hmm-filter}
\end{equation}
by the standard forward recursion, initialized from the stationary state
probabilities at the left edge of the window \citep{BaumPetrie1966}.  The target
is the posterior regime score
\begin{equation}
    F_{\mathrm{HMM}}(\bm X_{t,w})
    :=\sum_{j=1}^3 a_j\pi_{t,w}(j).
    \label{eq:functional-hmm-score}
\end{equation}
The learner receives only the emission window; neither the latent state nor the
filtering probabilities are included in its input.

\paragraph{GARCH conditional-volatility state.}
For the GARCH process in \cref{app:garch-family}, let
\(
  \bar\sigma^2
  :=\omega/(1-\alpha_{\mathrm G}-\beta_{\mathrm G})
\)
be the stationary second moment under the finite-variance parameterization
\citep{Bollerslev1986}.  The bounded target
\begin{equation}
    F_{\mathrm{GARCH},t}
    :=
    \tanh\!\left[
      c_\sigma
      \log\!\left(\frac{\sigma_{t+1}^2}{\bar\sigma^2}\right)
    \right],
    \qquad c_\sigma=\frac12,
    \label{eq:functional-garch-volatility}
\end{equation}
asks the learner to recover the next conditional-volatility state from the
bounded return-derived input window.  The latent variance is used only to
construct the synthetic label.

\subsection{High-dimensional observation constructions}
\label{app:high-dimensional-functionals}

The high-dimensional experiment changes the observation map while keeping the
three-dimensional teacher fixed.  Let \(L_t\in[-1,1]^3\) be the bounded latent
VARMA process.  For every ambient dimension, the exponential and Volterra
targets are computed from \(L_{t-w+1:t}\) using \(u_3\) and \(v_3\).  The
learner instead receives either the sparse or distributed observation
\(X_t^{(d)}\) defined in \cref{app:varma-family}.  In the sparse construction,
the informative coordinates are the first three and all remaining coordinates
are nuisance.  In the distributed construction, the informative subspace is
spread through \(A_d\), which is generated from the named JL-independent
embedding stream and stored with the dataset artifact.  Thus label range and
intrinsic task difficulty do not increase merely because \(d\) increases.

This design distinguishes two possible compression failures.  Sparse signal
may be diluted by indiscriminate projection when only a few coordinates matter;
distributed signal asks whether a global low-dimensional geometry survives the
same fixed encoded width.  No target direction is selected after inspecting
prediction results.

\subsection{Implementation contract for target generation}
\label{app:functional-implementation}

A target generator receives a window ending at time \(t\) and an immutable
context object.  Window-only functionals use only the window.  The genuine
forecast functional reads \(X_{t+1}\) from the context, the HMM score reads the
precomputed finite-window filtering probabilities, and the GARCH target reads
\(\sigma_{t+1}^2\).  These fields are teacher-only and are excluded from the
model input and preprocessing.  The context, target parameters, and teacher
arrays are checksum-validated and stored separately from feature artifacts.

The implementation must not label \(u^\top X_t\) as one-step forecasting.
Current-point projection may be retained as an optional smoke test under the
identifier \texttt{current-point-decoding}, but it is not an E1 forecasting
result.  Delayed-recall and cross-lag targets are deterministic window-only
objects, while all direction vectors are created once per dataset artifact and
reused for every example and method.

\subsection{Assumption compatibility}
\label{app:functional-assumptions}

\begin{lemma}[The controlled functional suite satisfies the bounded-label and
mixing requirements]
\label{lem:functional-compatibility}
For each process realization in \cref{app:dependent-families}, all targets in
\cref{eq:functional-forecast,eq:functional-exp,eq:functional-volterra,eq:functional-delay,eq:functional-cross,eq:functional-hmm-score,eq:functional-garch-volatility}
are measurable and bounded by one in absolute value.  When paired with the
corresponding input process, the resulting supervised process is stationary and
beta mixing under the conditions of
\cref{lem:varma-beta,lem:hmm-beta,lem:garch-beta}.
\end{lemma}

\begin{proof}
The \(\ell_1\)-normalization of the directions and
\(X_t\in[-1,1]^d\) give the bounds for forecasting, delayed recall, and the
cross-lag product.  Equation~\cref{eq:functional-Aw} and the denominators in
\cref{eq:functional-exp,eq:functional-volterra} give the same bound for the two
fading-memory targets.  The HMM score is a convex combination of values in
\([-1,1]\), and the GARCH target is bounded by the outer hyperbolic tangent.
Each common task is a finite-block function of the underlying process. The
forecast target uses one future observation and therefore invokes
\cref{lem:finite-block-transform} with \(\ell_+=1\); the remaining shared tasks
have \(\ell_+=0\). The HMM score is a measurable function of a finite emission
block. Although it is indexed by \(t+1\),
\(\sigma_{t+1}^2=\omega+\alpha_{\mathrm G}R_t^2+
\beta_{\mathrm G}\sigma_t^2\) is a measurable function of the current
stationary GARCH state. The conclusion follows from
\cref{lem:finite-block-transform} and the family-specific mixing lemmas. No
unbounded label noise is added in the theory-aligned controlled suite.
\end{proof}

\subsection{Task allocation across the empirical study}
\label{app:functional-allocation}

The full Cartesian product of generators and targets is intentionally avoided.
\Cref{tab:functional-allocation} assigns each functional to a specific
scientific question.

\begin{table}[H]
\centering
\small
\caption{Compact allocation of controlled functionals.  ``Shared'' means that
the same mathematical target is applied to each listed generator; the
family-specific targets use latent variables only to construct labels.}
\label{tab:functional-allocation}
\begin{tabularx}{\linewidth}{@{}lXX@{}}
\toprule
Experiment & Generators & Functionals \\
\midrule
Functional hierarchy and learning curves
  & VARMA, medium persistence
  & future projection, exponential memory, Volterra, delay curve, cross-lag \\
Memory and persistence study
  & VARMA at low/high persistence
  & delayed recall at \(\tau\in\{1,10,20\}\) under a contraction sweep \\
Cross-family robustness
  & VARMA, HMM, GARCH at low/medium/high persistence
  & shared exponential and Volterra tasks; VARMA future projection; HMM filtered score; GARCH volatility state \\
Ambient-dimension study
  & latent VARMA observed at \(d\in\{3,30,100,300,500\}\)
  & latent exponential and Volterra teachers under sparse/distributed observations \\
Resource and memory ablations
  & VARMA, medium persistence
  & Volterra and delayed recall at \(\tau=20\) \\
Finite-shot study
  & VARMA, medium persistence
  & exponential memory, Volterra, and delayed recall at \(\tau=10\) \\
Chaotic stress tests
  & Lorenz--63 and Mackey--Glass
  & multi-horizon forecasting and partial-observation prediction; see \cref{app:chaotic-stress} \\
Real-world sanity checks
  & Beijing PM2.5, live fuel moisture, hydraulic systems
  & archive-provided scalar regression targets; see \cref{app:real-world-sanity} \\
IBM hardware feasibility
  & compact \(d=2\) VARMA subset
  & delayed recall at \(\tau\in\{1,3\}\); see \cref{app:hardware-validation} \\
\bottomrule
\end{tabularx}
\end{table}
